\documentclass[11pt]{article}
\usepackage{amsmath,amssymb,hyperref}
\usepackage[margin=1in]{geometry}
\title{Beyond Linear Head Savings: Logarithmic-Saving Audit}
\author{}
\date{November 23, 2026}
\begin{document}
\maketitle

\paragraph{Outcome.}
The revised exact target
$W(n)=O(2^n/(n\log\log n))$ was not proved, either in Lean or in the
written arguments. The earlier positive-exponent target and both
fallbacks were also not reached: this run proves no bound
$W(n)=O(2^n/(n g(n)))$ with $g(n)\to\infty$, and no unrestricted
lower bound $W(n)\geq c2^n/n$ on infinitely many dimensions.
The new results concern necessary information, a specified class of
localized constructions, and exact failures of proposed transfer steps.
They do not close the exponent gap.

The instruction to assume an affirmative result exists was treated as
a research direction, not as an axiom or a premise of a claimed proof.
No restricted-domain construction below is asserted to meet the target.

\section{Verified baseline and unchanged scope}

The repository's existing theorem
\texttt{worstCaseHeadComplexity\_real\_le\_eleven\_mul\_two\_pow\_div}
in \texttt{Results/UniversalLinearSaving.lean} states, for $n>0$,
\[
 (\texttt{HeadComplexity.worstCaseHeadComplexity}\ n:\mathbb R)
 \leq 11\,2^n/n.
\]
The new modules import the unchanged definitions of the head resource,
the one-layer attention model, and admissibility. In particular,
positivity of an affine denominator on the cube alone is insufficient
to identify it with an atom: its nonzero slopes must have one common
orientation. Strict signs are used throughout.

The existing counting file has an important formal-status qualification.
\texttt{worstCaseHeadComplexity\_gt\_div\_of\_warren}
takes \texttt{WarrenStrictSignPatternBound} as a hypothesis.
The reduction to parameter polynomials is Lean-checked; Warren's
general sign-pattern theorem is an explicit interface, not a proved
theorem of that file. The familiar mathematical
$\Omega(2^n/n^2)$ conclusion follows by the classical theorem, but
should not be described as an unconditional Lean theorem here.
No axiom was added to discharge this interface.

The separately supplied branch
\texttt{ar/1dc3df1f7016} of
\texttt{VerifiedMechanisms/autoresearch-runs} was inspected at
\texttt{045a8fb}. Its declaration
\texttt{HeadBound.HStar\_le\_all}, in
\texttt{HeadBound/SubCube.lean}, bounds every function by
$2^n/2^{\lfloor\log_2 n\rfloor}$, using natural-number division.
The branch's compiler first proves full linear spanning and then
interpolates the truth table exactly. That mechanism does not transfer
a lower-dimensional sign-only local certificate.

A final cross-run check read 91 run reports at the 31 branch tips of
\texttt{autoresearch-runs} and 70 branch tips of
\texttt{xorformer-runs}; ten of the latter lacked a conventional
report. No claimed universal exponent improvement or superconstant
saving was found in those reports. Restricted quadratic-function
bounds and a last-token-readout variant do not meet this task.
This was a report audit, not an independent rebuild of every branch.

\section{New results, separated by proof status}

All Lean paths below are relative to
\texttt{formalization/HeadComplexity/Results/}; all names are in
namespace \texttt{HeadComplexity}.

\paragraph{Finite precision and denominator information, Lean-checked.}
\texttt{PrecisionBarrier.lean} proves
\texttt{two\_pow\_le\_bits\_mul\_affine\_parameters\_of\_finiteAlphabet\_universal}
and
\texttt{two\_pow\_le\_bits\_mul\_raw\_parameters\_of\_finiteAlphabet\_universal}.
For a fixed alphabet of at most $2^b$ elements these give, respectively,
\[
 2^n\leq b(1+2H(n+1)),\qquad
 2^n\leq b(1+H(4+2n)).
\]
The first chart encodes effective affine numerators and denominators;
the second encodes the original \texttt{FracAtom} scalar parameters.
The hypotheses require universality over all strict Boolean tables.

\texttt{FixedDenominatorLibrary.lean} proves a finite Sauer bound,
then the denominator-only result
\texttt{two\_pow\_two\_pow\_le\_bits\_denominator\_capacity}:
\[
 2^{2^n}\leq 2^{bnH}
       \sum_{j=0}^{1+nH}\binom{2^n}{j}.
\]
Here only the normalized denominator coefficients have a finite
alphabet. Every numerator coefficient and the bias may be any real
number. This result is unconditional and does not use Warren's theorem.

\paragraph{Asymptotic information consequences, written only.}
The complete proof in \texttt{exponent-precision.tex} shows that, for
every $g(n)\to\infty$ and fixed $C>0$, a universal construction with
$H\leq C2^n/(n g(n))$ requires a denominator library of size $K$ with
\[
 \log_2 K\geq(1-o(1))2^n.
\]
Thus normalized denominator coefficients in fixed $b$-bit alphabets
require $b\geq(1/C-o(1))g(n)$. Counting all effective coefficients
gives the weaker constant $1/(2C)$. The earlier precision claim is
therefore valid in these intrinsic finite-information interpretations,
with the head-count constant included. It is a worst-case information
bound, not a requirement on each individual weight of each target.

\paragraph{Raw width qualification, partly Lean-checked.}
The actual model leaves width unrestricted. The written construction
in \texttt{exponent-precision.tex} realizes every strict $H$-head
certificate, for $H\geq1$, with the same number of heads, larger finite
width, shared embeddings, and all raw entries in
$\{0,1,-1,\sqrt2,-\sqrt2\}$. Integer multiplicities are carried by
width; this is not a finite binary-rational precision theorem.
The construction uses positivity of exponential denominators and a
finite strict margin, so it preserves admissibility and every sign.

The approximation lemmas are formalized in
\texttt{RawAlphabetSimulation.lean}:
\texttt{dense\_rawAlphabetSum},
\texttt{dense\_rawAlphabetSum\_vectors}, and
\texttt{exists\_rawAlphabetSum\_strict\_certificate}.
The full embedding and matrix construction remains written only.

\paragraph{Adaptive local masks, Lean key lemma and written consequence.}
\texttt{ExponentMechanisms.lean} proves
\texttt{activeCountingClearedPoly\_totalDegree\_le},
\texttt{eval\_activeCountingClearedPoly}, and
\texttt{strictPolynomialSignPattern\_of\_active\_score}.
If at most $L$ heads are active at any input, pointwise denominator
clearing has parameter degree at most $L+1$.
\texttt{active\_strictPatterns\_card\_le\_warren} gives the
corresponding count conditional on the existing Warren interface.

Using the classical theorem, \texttt{exponent-mechanisms.tex} proves:
if the mask dictionary has size at most $2^{bn}$ and a universal
masked construction uses $H\leq C2^n/(ng(n))$, then
\[
 \liminf_{n\to\infty}
 \frac{\log_2(L_n+1)}{g(n)}\geq\frac1{2C}.
\]
This rules out bounded active-head overlap for every superconstant
saving, and polynomial active-head overlap for every positive power
saving. It counts heads, not cells. In particular, a cell carrying
$O(n/\log n)$ heads is not excluded from a logarithmic-saving route.
The theorem is about explicitly masked scores, or actual scores with
a proved sign-preserving truncation. It places no such restriction
on arbitrary attention representations with large cancellations.

\paragraph{A new partial ranking mask and its limitation, Lean-checked algebra.}
\texttt{PartialRankingMask.lean} proves
\texttt{partialRankingMask\_eq}, the identity
\[
 \frac{u-v}{p+u+v}-\frac{u}{p+u}+\frac{v}{p+v}
 =\frac{(u-v)uv}{(p+u+v)(p+u)(p+v)}.
\]
It cancels exactly when either local block is absent.
\texttt{partialRankingMask\_strictAnti\_parameter} and
\texttt{not\_partialRankingMask\_scalar\_repair} certify a concrete
failure of scalar layer repair. The actual three-head realization,
the uniform ranking approximation, and the six-bit instance are
written in \texttt{exponent-literature.tex}. The formal algebra is
not presented as a formalization of that complete compiler.

The same module also proves \texttt{reversePartialRankingMask\_eq}
and \texttt{reversePartialRankingMask\_pos}. Replacing each denominator
by $p-u-v$, $p-u$, or $p-v$, with $p>u+v$, gives the analogous
axis-zero identity in the orientation needed by the gluing theorem
below. Its comparison amplitudes are not uniformly saturated.

\paragraph{Exact localization, Lean singleton theorem and written extension.}
\texttt{ExactMaskDegree.lean} proves
\texttt{le\_totalDegree\_of\_cube\_supported\_at\_zero},
\texttt{le\_heads\_of\_clearedAffinePoly\_supported\_at\_zero}, and
\texttt{le\_heads\_of\_atom\_score\_supported\_at\_zero}.
In particular, a nonzero score of actual admissible atoms that is
exactly zero at all cube vertices except the all-false vertex needs
at least $n$ heads. This is an exact-zero support statement, not a
lower bound for sign computation of the singleton indicator. The
radius-$r$ extension by restriction to a maximal support vertex is
written in \texttt{exponent-literature.tex}.

\paragraph{A sign-only interval decoder, written construction and Lean kernels.}
\texttt{torus-prism-decoder.tex} gives a complete extraction construction:
for all sufficiently large $m$, every $m$ by $s$ Boolean table on
the two-one-hot inputs has an $O(m/\log m)$-head representation,
uniformly in $s$. All denominators are admissible on the entire
$(m+s)$-cube, not merely on the promised inputs. The score has a
strict uniform margin on the grid. A finite bias perturbation makes
every other output strict without prescribing its label.
The proof uses a finite Hall-matching extraction array and differences
of rational suffix kernels. A related zero-bias decoder vanishes
exactly on both coordinate axes, is nonnegative everywhere, and
approximates every binary or ternary grid table with arbitrary
specified error at the same head order. Its denominator orientation
is opposite to the one required by the gluing theorem below.

\texttt{TorusPrismDecoder.lean} proves the nine public kernel lemmas,
including \texttt{rationalPrefixKernel\_error},
\texttt{rationalSuffixKernel\_error},
\texttt{gatedRationalInterval\_eq},
\texttt{gatedRationalInterval\_nonneg}, and the two exact zero lemmas.
The combinatorial extraction construction and its complete compiler
are written only. The note also computes the decoder's limit on every
cube input, so its uncontrolled labels outside the grid are explicit.

\paragraph{Complementary rails, Lean-checked exact transfer.}
\texttt{DualRailLift.lean} proves seven public theorems. In particular,
\texttt{exists\_fracAtom\_dualRail\_affineRatio} realizes every logical
affine ratio with positive denominator on an $n$-cube by one actual
\texttt{FracAtom} on $2n$ physical coordinates, exactly on inputs
$(x,1-x)$. The lifted denominator is positive throughout the entire
physical cube and has strictly positive slopes. The proof includes
the exact positivity margin and all intermediate identities.
This is not a compiler on the original $n$-cube, nor a representation
of arbitrary labels on the entire $2n$-cube. No head bound for $W$
is inferred by changing dimensions or restricting the input promise.

\paragraph{Periodic-chain orientation repair, written only.}
\texttt{exponent-literature.tex}, Proposition~\texttt{prop:periodic-chain},
proves that an $H$-term arithmetic pole chain with $k$ negative slopes,
sinusoidal residues of period $k+1$, and the specified positivity
buffer can be rewritten exactly using at most
$H+k(k+1)(k+2)/3$ ratios with all slopes positive.
\texttt{exponent-mechanisms.tex}, Theorem~\texttt{thm:hook-repair},
extends this via explicit Laurent interpolation to zero-mean periodic
affine numerators, with $H+O(k^2L+k^3)$ heads and buffer
$\min D>\delta[kL+k(k-1)/2]$, where $b_i<L\delta$.
These are exact whole-cube identities and therefore preserve any
existing strict signs. Their periodic residue and arithmetic pole
requirements are substantive.

\texttt{periodic-chain-decoding.tex} proves that repeating a complete
zero-mean period of length $K$ produces at most $K-2$ scalar zeros,
regardless of the number of periods. Even when residues, shift, and
step depend jointly and affinely on one memory scalar, the zero count
is at most $4K-1$. Thus that decoder can shatter only $O(\log K)$
scalar probes by this mechanism. These bounds do not constrain a
general system with several independent nonlinear memory parameters.

\paragraph{Further explicit failed transfers, written only.}
The mixed-priority-list theorem in \texttt{exponent-mechanisms.tex}
proves a worst-case $\Omega(2^n/n)$ lower bound for sums of real
priority lists, allowing both orientations. Each list is sparse in
the fixed dictionary of positive and negative literal monomials;
counting dictionary supports and linear sign patterns proves the
bound. This is a subclass lower bound, not item 3 of the original task.
The same note gives a grid-routing local compiler, exact three-head
single-slope reversal, and the support cost of weak affine-flat
certificates. None yields an unrestricted asymptotic saving.
\texttt{product-tangent-counterexample.tex} supplies explicit positive
affine factors for which no common interpolation parameter makes
the logarithmic derivative agree in sign with the product comparison
on even the two inputs of a one-bit cube. Pointwise mean-value
arguments therefore do not supply the proposed common tangent.

\paragraph{A valid global composition theorem, written only.}
The finite gluing theorem in \texttt{exponent-mechanisms.tex}
composes product downward-star modules without extra heads, when
their denominators have positive slopes in the original coordinates
and their scores vanish exactly on both local axes. The proof orders
centers by cardinality, rescales each module to correct its cell,
and adds a large positive outside-coordinate penalty to its
denominators to preserve all previously fixed strict signs. It
verifies admissibility on the whole cube. This proves a transfer
step but does not supply the required efficient local modules.

\paragraph{Opposite-orientation vote transfer, written only.}
For an $m$ by $m$ interleaved-order comparator, suppose $H$ heads
whose denominator slopes have the opposite orientation approximate
every numerical vote in $\{-1,1\}$ to error $\varepsilon$.
The complete proof in \texttt{exponent-literature.tex} gives
\[
 H2^H\geq m/2-(m\varepsilon+4H)^2
 \qquad\text{when }m\geq8H.
\]
It allows arbitrary numerator sizes and cancellation. In particular,
error $O(\log m/m)$ requires
$H\geq\log_2m-\log_2\log_2m-O(1)$.
The direct uniform-error summation guarantee for
$O(m/\log m)$ majority votes with integer margin one uses that
accuracy. Thus this
particular reversal of the known local compiler loses its saving.
This is a numerical-approximation obstruction, not a lower bound
for sign-only computation of that comparator, which is affine.

\section{Routes tried and the remaining lemma}

The literature transfer audit is in \texttt{exponent-literature.tex}.
It includes depth-two thresholds, deep memorization, sign rank,
rational networks, covering codes, and number-theoretic encodings.
Every successful transfer is stated on its exact domain. Both the
Erd\H{o}s--Moser majority-order representation and the new extraction
decoder give actual admissible $O(n/\log n)$-head sign representations
of arbitrary two-one-hot tables on $\Theta(n^2)$ inputs. Neither
prescribes arbitrary strict labels on the rest of the cube.

The failed extensions have explicit endpoints. Exact zero localization
to a radius-$r$ Hamming ball requires at least $n-r$ heads. Sums of
common-orientation first-hit decision lists representing parity need
at least $\binom n{\lfloor n/2\rfloor}$ lists. The new three-ratio
mask leaves a cylinder of unwanted signs; its proposed scalar repair
fails on an explicit pair of cube inputs. The direct downset gate
would also introduce inadmissible mixed denominator slopes. These
statements exclude the specified constructions, not arbitrary globally
cancelling head systems.

\paragraph{Single missing lemma.}
The concrete remaining lemma is the \emph{reverse-orientation
axis-zero local compiler}. For every $r,s\geq1$ and every strict
$r$ by $s$ sign table, construct a sum of at most
\[
 K\left(1+\frac{r+s}{\log\log(r+s+4)}\right)
\]
affine ratios with strictly negative denominator slopes in the deletion
variables, positive denominators on the entire $(r+s)$-cube,
the prescribed strict signs on $(e_i,e_j)$, and exact zero whenever
either variable block is all zero. All coefficients may depend on
the table. This lemma is unproved.

The written gluing theorem would apply after changing deletion
variables to original coordinates. Existing star covers provide
$O(2^n/n^2)$ cells, the local cost is uniformly $O(n/\log\log n)$,
and the two exceptional faces cost only
$O(2^{\lceil n/2\rceil})$ additional heads. Therefore this one lemma
would establish the revised exact target
$W(n)=O(2^n/(n\log\log n))$. Full quantifiers and the face
correction are in the last section of
\texttt{exponent-literature.tex}. The lemma remains a substantial
unresolved local universality problem; this implication is not
counted as a successful asymptotic result.

\section{Changes and verification}

No original definition or theorem statement was weakened or altered.
New auxiliary predicates, parameterizations, and score functions were
added in the new modules. \texttt{WorstCaseExponentImprovement} and
\texttt{WorstCaseLogLogSaving} in \texttt{PrecisionBarrier.lean} are
only proposition definitions of the respective eventual bounds
against \texttt{worstCaseHeadComplexity}; neither is asserted as a
theorem. The latter uses $n\log(\log n)$ with natural logarithms,
positive existential constant, and an arbitrary eventual threshold.
The results aggregate imports the
new modules. The previous \texttt{mechanism-audit.tex} and
\texttt{scientist-report.tex} are unchanged.

The complete \texttt{HeadComplexity} target was built with
\texttt{lake}. Axiom audits of all 53 new public proved declarations use only
\texttt{propext}, \texttt{Classical.choice}, and \texttt{Quot.sound}.
The Warren consequence retains its explicit theorem hypothesis.
There are no new axioms or admitted proofs. TeX sources were checked
for balanced environments and braces; no TeX compiler was available
on the machine, so PDF rendering was not checked.
The machine-readable file \texttt{exponent-validation.json} records
the checked source hashes and declaration names.

\end{document}
